\scanpage{69}
$\Delta=(S,A,T)$ graphe cubique

$\unsure{\widehat X|\chi|X}:X\mathrel{\underset{0}{\subset}}S$ pavé (anciennement : bitriangle)

1) $\forall x\in S-X,\ \mathrm{Omb}_X(x)\deleted{\unsure{X}}=\{s\in X\mid s\ \text{lié à }x\}$ est une triade de $X$. L'application
\[
\varphi_X:x\longmapsto\mathrm{Om}_X(x):\quad S-X\longrightarrow\mathrm{Tria}(X)
\]
est une application surjective, \deleted{dont les fibres sont des triades.}

2) Soit $\underline t=\{\alpha,\beta,\gamma\}$ une triade de $X$, $p_X^{-1}(\{\underline t\})$ l'ensemble des $x\in S-X$ d'ombre $\underline t$. Alors $\overline{\underline t}={\unsure{\tau_x|p_x|\varphi_X}}^{-1}\{(\underline t)\}\inserted{\unsure{\tau}}$ est une triade \deleted{notée $\psi(\underline t)$}, c'est l'unique triade \inserted{$\overline\tau$} de $S$ ``opposée'' à $\underline t$ i.e \inserted{disjointe de $\underline t$ et} dont les él. sont \deleted{\unsure{non}} liés à tous les éléments de $\underline t$ ; c'est aussi l'ens des $s\in S-\underline t$ \deleted{\unsure{non}} liés aux éléments de $\underline t$. \inserted{Dnc $\underline t$ est l'unique triade opposée à $\tau$.}

3) Les deux triades $\underline t',\underline t''$ de \deleted{\unsure{$S$}} $X$ qui sont parallèles à $\underline t$ \inserted{(i.e. ne rencontrant pas $\underline t$)} sont les deux uniques triades de $S$ \deleted{disjointes} \inserted{``liées'' à $\underline t$} (i.e. dont les él. sont \uline{non} liés à tous les éléments de $\underline t$) -- $\underline t'\cup\underline t''$ est l'ens des $s\in S-\tau$ qui sont non liés aux él. de $\tau$.

4) Les \deleted{réunion des $\psi(\underline t)\cup\psi(\underline t')\cup\psi(\underline t'')$ est un} \inserted{triades $\overline\tau,\overline\tau',\overline\tau''$ sont les triades parallèles} d'un pavé \inserted{$X_2$}, et les 6 triades $\tau,\tau',\tau''$ et $\overline\tau,\overline\tau',\overline\tau''$ forment un hexagone de triades liées

\diagram{\begin{tikzcd}[ampersand replacement=\&,row sep=tiny,column sep=small]
\&\overline\tau\arrow[dl,dash]\arrow[dr,dash]\arrow[ddl,dashed,no head]\arrow[ddr,dashed,no head]\&\\
\tau''\arrow[d,dash]\arrow[rr,dashed,no head]\arrow[ddr,dashed,no head]\&\&\tau'\arrow[d,dash]\arrow[ddl,dashed,no head]\\
\overline\tau'\arrow[dr,dash]\arrow[rr,dashed,no head]\&\&\overline\tau''\arrow[dl,dash]\\
\&\tau\&
\end{tikzcd}}

(on a indiqué en pointillé les paires de triades ``déliées'' \inserted{ou ``parallèles''}

% archive sheet number: 68

